\documentclass[11pt, a4paper]{article}

\usepackage[T1]{fontenc}
\usepackage[utf8]{inputenc}
\usepackage{tgpagella}
\usepackage[spanish, es-noshorthands]{babel}
\usepackage{amsmath, amssymb, bm}
\usepackage{booktabs}
\usepackage{tabularx}
\usepackage[table, xcdraw]{xcolor}
\usepackage[most]{tcolorbox}
\usepackage[american]{circuitikz}
\usepackage[margin=2.5cm]{geometry}
\usepackage{fancyhdr}

\pagestyle{fancy}
\fancyhf{}
\setlength{\headheight}{16pt}
\lhead{\textbf{UCB Sede Tarija} | Ing. Mecatrónica}
\chead{}
\rhead{IMT-121 Circuitos Electrónicos I | Ejercicios S06}
\lfoot{Prof: M.Sc. Ing. Bernardo Quiroga Turdera}
\rfoot{Página \thepage}
\renewcommand{\headrulewidth}{0.4pt}

\definecolor{ucbblue}{RGB}{0, 51, 102}

\begin{document}

\begin{tcolorbox}[colback=ucbblue!5!white, colframe=ucbblue, title=\textbf{Hoja de Ejercicios: Teorema de Superposición}]
    \centering \large \textbf{Análisis de Redes Resistivas con Múltiples Fuentes}
\end{tcolorbox}

\section*{Ejercicio 1: Aplicación Completa del Método}
Considera el circuito alimentado por una fuente de tensión y una fuente de corriente:
\begin{itemize}
    \item $V_s = 9\,\text{V}$ (terminal positivo arriba, negativo a tierra).
    \item $R_1 = 3\,\text{k}\Omega$ (conectado entre el positivo de $V_s$ y el nodo de salida $V_o$).
    \item $R_2 = 6\,\text{k}\Omega$ (conectado entre el nodo $V_o$ y tierra).
    \item $I_s = 1\,\text{mA}$ (conectado en el nodo $V_o$, extrayendo corriente hacia tierra).
    \item \textbf{Variable de interés:} Tensión $V_o$ respecto a tierra.
\end{itemize}

\begin{enumerate}
    \item Dibuja el subcircuito equivalente manteniendo activa únicamente la fuente de tensión $V_s$ (desactiva $I_s$). Calcula la contribución $V_o^{(1)}$.
    \item Dibuja el subcircuito equivalente manteniendo activa únicamente la fuente de corriente $I_s$ (desactiva $V_s$). Calcula la contribución $V_o^{(2)}$. \textit{Atención al sentido de la corriente y la referencia de $V_o$.}
    \item Realiza la suma algebraica para obtener la tensión total $V_o$.
    \item Verifica tu resultado planteando el análisis nodal (KCL) en el circuito original.
    \item Explica por qué una de las contribuciones resulta negativa.
\end{enumerate}

\section*{Ejercicio 2: Diagnóstico de Errores}
Un estudiante analiza un circuito y calcula los siguientes valores parciales para la tensión de salida ante dos fuentes distintas:
\begin{equation*}
    V_o^{(1)} = 6\,\text{V}, \quad V_o^{(2)} = 2\,\text{V}
\end{equation*}
El estudiante concluye inmediatamente que $V_o = 8\,\text{V}$. 
\begin{enumerate}
    \item ¿Qué aspecto fundamental del análisis por superposición omitió el estudiante?
    \item ¿Cómo influyen las referencias de polaridad y las direcciones de corriente en el cálculo final?
\end{enumerate}

\section*{Ejercicio 3: Razonamiento Físico de Desactivación}
Explica con tus propias palabras, basándote en las Leyes de Kirchhoff y las definiciones de fuentes ideales:
\begin{enumerate}
    \item ¿Por qué desactivar una fuente de tensión ideal ($V_s = 0\,\text{V}$) equivale a un cortocircuito?
    \item ¿Por qué desactivar una fuente de corriente ideal ($I_s = 0\,\text{A}$) equivale a un circuito abierto?
\end{enumerate}

\section*{Ejercicio 4: Limitación de Potencia}
En un circuito lineal, se determinó mediante superposición que la tensión sobre un resistor de $R = 2\,\text{k}\Omega$ debido a dos fuentes separadas es $v_1 = 3\,\text{V}$ y $v_2 = -1\,\text{V}$.
\begin{enumerate}
    \item Calcula la tensión total $v_{\text{total}}$ sobre el resistor.
    \item Calcula la potencia total disipada usando $v_{\text{total}}$.
    \item Calcula las potencias parciales $P_1 = v_1^2/R$ y $P_2 = v_2^2/R$, y comprueba si $P_1 + P_2$ es igual a la potencia total. Explica la discrepancia matemáticamente.
\end{enumerate}

\end{document}
